\scanpage{69}
\mspage{34}
et de reconstituer en principe les schémas lisses sur des $K$, par recollements de tels morceaux avec des ``imm. ouvertes''. Mais je sens nettement que pour faire une telle description, il me faudra développer un langage géométrique qui colle mieux à l'intuition géométrique, que les sempiternelles extensions de groupes profinis\ldots\ ou actions extérieures, et où les points rigides (à coeff dans des clôtures algébriques de corps) jouent un rôle prépondérant. Il \deleted{me} faudra y revenir dessus -- et en même temps, expliciter les topos étales (pas seulement le ``morceau $K(\pi,1)$) tout entiers des schémas décrits ici par des extensions.
\margin{68}
